% Paper 3 — Related Work Section
% Building Local-Language Economic Narrative Indices

\section{Related Work}\label{sec:related}

\subsection{Text-as-Data in Economics}

The use of textual data in economic research has evolved through several methodological eras. Early work by \citet{tetlock2007giving} demonstrated that the pessimistic tone of Wall Street Journal columns predicted downward pressure on market returns, using the General Inquirer dictionary for sentiment extraction. \citet{tetlock2007sad} extended this approach to individual stock predictions, showing that negative media sentiment predicted lower firm earnings and higher market returns.

The Economic Policy Uncertainty (EPU) index of \citet{bakereconomic} established the template for news-based economic indices: count articles containing specific keyword combinations related to the economy, policy uncertainty, and future-oriented language. The EPU index has become one of the most widely cited text-as-data instruments in economics, spawning country-specific variants for China \citep{huang2018measuring}, Europe \citep{ferrara2022uncertainty}, and other regions.

More recent work has moved beyond dictionaries to machine learning. \citet{thorsrud2018text} applied structural topic models to Norwegian news for nowcasting GDP. \citet{jurado2015macroeconomic} used text-derived uncertainty measures in macroeconomic forecasting. \citet{nabil2026economic} conducted a systematic review of 66 empirical studies on sentiment-based economic forecasting, finding a median RMSE improvement of 10--15\% after correcting for publication bias.

Despite this growth, the field remains heavily concentrated. \citet{nabil2026economic} document that 84\% of studies use English-language text, 56\% focus on the United States, and no sentiment indices exist for Bangla, Swahili, Hindi, or any other South Asian or African language.

\subsection{Local-Language Economic Indices}

The construction of economic text indices in languages other than English has been limited. Chinese-language EPU indices have been developed for mainland China \citep{huang2018measuring} and Taiwan. Japanese sentiment indices exist for financial markets. Arabic-language sentiment analysis has been applied to Gulf economies.

However, these efforts share a common feature: they target countries with developed financial markets, strong institutional data infrastructure, and substantial academic communities producing English-language publications. For low-income, high-population economies---where official statistics arrive late, informal activity is large, and local-language news is the primary information channel---no comparable indices exist.

The gap is particularly acute for Bangla. Despite Bangladesh's 2025 economy of roughly \$460 billion, its growing garment export sector, and its frequent appearance in development economics literature, no reproducible text-based economic indicator exists in Bangla. This paper addresses that measurement gap by documenting the BENI v1 database and method pipeline.

\subsection{Low-Resource NLP for Economic Text}

NLP in low-resource languages faces well-documented challenges: limited labeled data, underperforming tokenizers, morphological complexity, and code-mixing between Bangla and English in online text \citep{devlin2019bert}. For economic text specifically, the challenge is compounded by domain-specific vocabulary (monetary policy terms, trade jargon, fiscal terminology) that general-purpose language models handle poorly.

BanglaBERT \citep{chowdhury2020banglabert} provides a pre-trained Electra-based model for Bangla that has achieved strong performance on classification tasks. However, no prior work has fine-tuned BanglaBERT for economic text classification or applied it to construct a narrative index.

\subsection{Annotation Cost in Low-Resource Settings}

A critical but under-studied question is how much human annotation is needed to build a reliable classifier in a low-resource language. \citet{settles2009active} established the theoretical foundations of active learning, showing that uncertainty sampling can achieve comparable performance to random sampling with 30--50\% fewer labeled examples. Recent work in NLP has applied active learning to sentiment analysis \citep{shen2018deep}, named entity recognition, and other tasks.

However, no prior work has measured the annotation cost curve for economic text classification in a low-resource language. This gap matters because annotation is the binding constraint for researchers working outside well-funded NLP laboratories. If 150 carefully selected articles can achieve 95\% of the performance of 300 articles, the practical barrier to building local-language economic indices drops substantially.

\subsection{Our Contribution}

This paper makes four contributions to the literature:

\begin{enumerate}
    \item \textbf{Legacy prototype benchmark}: We compare TF-IDF and BanglaBERT on a silver-standard dataset of 299 LLM-annotated articles. The central finding is that label quality dominates model architecture: models trained on the same weak labels produce statistically indistinguishable errors (McNemar $\chi^2 = 0.18$, $p = 0.67$), and the simpler TF-IDF model outperforms BanglaBERT (F1 = 0.44 vs. 0.16). The LLM labels are treated as reference labels, not as independent human gold-standard labels.
    \item \textbf{Annotation cost curve}: We use active learning simulation to measure how classification performance scales with annotation investment. At 7.4\% base rate, TF-IDF cannot learn the minority class at any budget up to $k = 299$, establishing a lower bound on the annotation investment needed for rare-event economic text classification.
    \item \textbf{Replicable pipeline}: We provide a complete, open-source pipeline from raw newspaper articles to validated narrative index, designed for replication in any low-resource language.
    \item \textbf{Local-language advantage}: We demonstrate that Bangla news captures long-run macroeconomic co-movement, establishing the viability of local-language text as economic data.
\end{enumerate}
